%%%%%%%%%%%%%%%%%%%%%%%%%%
% LSHTM Med Stats MSc
% Project Template
% Instructions:
% - Compile only main.tex (this file)
% - Make formating choices here
% - Put content in seperate .tex files and input them here
% - See chapter02.tex for examples of things
%%%%%%%%%%%%%%%%%%%%%%%%%%

\documentclass[11pt,a4paper,twoside,openany]{book}
% twoside: slight warning if printed will give unequal margins to allow for binding
% can be adjusted if needed; openany can be changed to open right if you want chapters
% to start on the right

% Core Formating Packages:
\usepackage{array,       % Lets you adjust table spacing
            ragged2e,    % For some layout issues 
            graphicx,    % 
            amsmath,     % 
            geometry,    % Page layout/Margins
            setspace,    % 1.5 spacing
            fancyhdr,    % Header editing
            hyperref,    % Allows links
            %indentfirst, % Uncomment for 1st para indentation
            titlesec     % Formating of chapter titles
            }

%%%% ADDITIONAL PACKAGES: %%%%          
\usepackage{helvet}
% Arial is the recommended font for MSc projects at LSHTM (bleh)
\renewcommand{\familydefault}{\sfdefault} % Sets default font to helvetica
% Arial is recommended font, Helvetica is almost identical
% If you must use Arial can use it with fontspec pkg and
% XeLatex (bleh).

\usepackage[justification=justified,   % Justifies captions
            singlelinecheck=false,     % Multiline captions with //
            labelsep=period,           % . sep, ex: Figure 1. instead of Figure 1:
            labelfont=bf]{caption}     % Bold figure/table title

\usepackage[backend=bibtex]{biblatex}
\bibliography{references.bib} % Your Bibtex file goes here
% Depending on what program you use to compile, you may
% want to remove the [backend=bibtex]

%%%% FORMATTING: %%%%
% Header:
\pagestyle{fancy}
\fancyhead{}
\fancyhead[RO]{\footnotesize\rightmark}   % Section title on right, odd pages
\fancyhead[LE]{\footnotesize\leftmark}    % Chapter Title on left, even pages
\fancyfoot[CE,CO]{\thepage}               % Centered page # in footer
\renewcommand{\headrulewidth}{0.0pt}      % Line below header
\renewcommand{\footrulewidth}{0.0pt}      % Line above footer

% Link & Metadata:
\hypersetup{
  colorlinks = true,
  urlcolor = black,    % External Links, default: blue                
  linkcolor = black,   % Internal Links, default: red  
  citecolor= black,    % citations, deafault: lime green  (?)
  pdfauthor = {},      % Don't put your name or it will show up in metadata
  pdfkeywords = {},
  pdftitle = {MSc Project Report},
  pdfsubject = {},
  pdfpagemode = UseNone
}
 
% Path for Images:
\graphicspath{{figures/}}  

% Margins:
\geometry{
  left=1in, right=1in,
  top=1in, bottom=1in,
  head=0.5in, foot=0.5in
}

% Chapter Title Formating w/ titlesec pkg:
\titleformat{\chapter}{\normalfont\huge}{\thechapter.}{20pt}{\huge\bf\scshape}
% scshape: Big/Small Caps
% Removes "Chapter X" titles in header & before text:
\renewcommand{\chaptername}{}

% Prevents latex from vertically trying to 
% fill pages with ragged2e command:
\raggedbottom

%%%% Other options you may want %%%%
% Used to change table horzontal spacing:
% \renewcommand{\arraystretch}{1.0}
% Table vertical spacing:
 {\setlength{\extrarowheight}{1pt}
% Used to change default indent spacing:
% \setlength{\parindent}{2em}
% Spacing between paragraphs:
% \setlength{\parskip}{0em}

%%%%%%%%%%%%%%%%%%%%%%%%%%%
% DOCUMENT
%%%%%%%%%%%%%%%%%%%%%%%%%%%

\begin{document}

%%%% Front Matter %%%%
\frontmatter

\begin{titlepage}

  \begin{flushright}
   \includegraphics[width=0.3\textwidth]{schoolseal.jpg}
  \end{flushright}
  
  \begin{center}
    \LARGE
        MSc Project Report \\
        2017--2018
    \vspace*{1cm}
        
    \Huge
        \textbf{Title of the Report}
    \vspace{2.5cm}
        
    \LARGE
        \textbf{Candidate Number:} 12345678 \\
        \textbf{Supervisor:} Dr. Supervisor \\
    \vspace{0.5cm}
    \large
        Page Count:
    \vfill
    \Large
        Submitted in part fulfilment of the requirements for \\
        the degree of MSc in Medical Statistics \\
    \vspace{0.8cm}
    \large
        September 2017 \\
        
  \end{center}
  
\end{titlepage}

  % Title Page

\onehalfspacing         % 1.5 Spacing, for double: \doublespacing
  \begin{center}
\textbf{Acknowledgements}
\end{center}
\addcontentsline{toc}{section}{Acknowledgements}

Blah Blah
\newpage                % Seperate ack & abstract
  \begin{center}
\textbf{Abstract}
\end{center}
\addcontentsline{toc}{section}{Abstract}

Blah Blah

\tableofcontents
\listoffigures
\listoftables

%%%% Main Matter %%%%

\mainmatter
  \chapter{Introduction}

\section{Introduction Section 1}

\subsection{Introduction Subsection 1}
This is text for your introduction.
\subsection{Introduction Subsection 2}
\subsection{Introduction Subsection 3}


\section{Introduction Section 2}

\subsection{Introduction Subsection 1}
\subsection{Introduction Subsection 2}
\subsection{Introduction Subsection 3}

  \chapter{Which Ice Cream is the Best?}

\section{Clearly Rocky Road}
In this paragraph we will discuss the eternal debate surrounding the suppremacy of Rocky Road as the best icream flavor. \\

Who doesn't love how the creamy chocolate ice cream contrasts with both the marshmellows and almonds. \footnote{My coffee this morning is pretty terrible so I am not responsible for anything written here.}

\subsection{What is the ideal ratio of marshmellows to almonds?}

\newpage

This new page demonstrates the header. Look up there! \\

But while we are here, lets again mention that Rocky Road is the best flavor of icecream. Here is a citation that has nothing to do with ice cream: \cite{Yman2016}

\newpage

Hey look at the header on an opposite page. Look up! The indentations on this page are a little odd looking because we manually broke the pages up. Don't do that. \\

Coffee, Mint Chip, Napolean \& Pistachio are also good flavors. Look at the pretty pic in \ref{fig:logo}

\begin{figure} 
  \includegraphics[width=0.50\textwidth]{schoolseal.jpg} % Can adjust size <-
  \caption{This is the logo of the London School of Hygiene \& Medicine}
  \caption*{It has nothing to do with ice cream (sadly, am I right?). 
            This part of the caption won't show up in the TOC because of the *
            }
% If you want caption above, simply put it above "\includegraphics"
  \label{fig:logo}
\end{figure}

\begin{table}[bth] % [bth] Tells LaTeX we want the table at the bottom (b) if possible
                   % if not then the top (t) is next best, lastly here (h)
                   % if you put [h] alone, then God smites a puppy. Lastly [p] will put the 
                   % on a special page for just figures.
\small             % lets make our table small!
\caption{The best ice cream toppings ranked.}
\begin{tabular}{|lcc|} \firsthline  % The | lines here tell Latex we want vertical lines there
 Rank & Topping & Qualities \\ \cline{1-3}
 1 (the best) & Rainbow Sprinkles & They are amazing. Really there is nothing better. \\
 2-500 & Everything Else & Never as good as rainbow sprinkles. \\ \hline
\end{tabular}
\end{table}


  \chapter{Chapter 3}
  \chapter{Chapter 4}
  \chapter{Conclusion}

%%%% Back Matter %%%%

\backmatter
\printbibliography[title = {References}]
\addcontentsline{toc}{chapter}{References}
\appendix
\chapter{Appendix}

\end{document}